\documentclass{article}
\usepackage{../math_template}

\author{Jonathan Kasongo}
\title{USA Math Olympiad 1976 --- P5/5}

\begin{document}
\maketitle

\begin{problem}{USA Math Olympiad 1976 --- P5/5}
The polynomials $A(x), B(x), C(x)$ and $D(x)$ satisfy the equation

$$
A(x^5) + xB(x^5) + x^2 C(x^5) = (1+x+x^2+x^3+x^4)D(x)
$$

Show that $A(1)=0$.
\end{problem}

\begin{solution}{Solution}
Observe the following

\[
\begin{aligned}
A(x^5) + xB(x^5) + x^2 C(x^5) &= \frac{x^5 - 1}{x-1}D(x) \\
(x-1)[A(x^5) + xB(x^5) + x^2 C(x^5)] &= (x^5 - 1)D(x) \\
\end{aligned}
\]

Take a 5th root of unity $z$ that isn't 1, so $z = e^{2\pi i k/5}$ for
$k=1,2,3,4$. Then

$$
p(z) := A(1) + zB(1) + z^2C(1) = 0
$$

has $4 > \deg p = 2$ roots which implies that $p(z) = 0$ for all
$z\in \mathbb{C}$. Then since $p(0) = 0$ we have $A(1) = 0$. $\blacksquare$
\end{solution}

\end{document}
